
\chapter*{Abstract}
\addcontentsline{toc}{chapter}{Abstract}

Portfolio optimization pipelines usually forecast returns and then optimize over those forecasts, which trains the model on prediction error rather than the quality of the decisions it makes. This dissertation studies an alternative in which training is done end-to-end on the risk-adjusted returns of the realized stream, gated by regimes discovered during training rather than statistically. The base underlying model is a Simple Cycle Reservoir, which is deterministically obtained and frozen, making the entire pipeline lightweight in both training and inference, making it cheap to scale. The central research question is whether such a simple reservoir model can outperform other popular methods like 1/N, Buy and hold SPY, Ledoit Wolf MV, LSTM, etc and whether decision-defined regimes carve the market differently than statistical regimes, and, if so, what underlying dimensions they vary upon.

I find that overall my methodology performs similarly to or worse than those baselines on the Sharpe ratio. However, I find that the learned gates carve the market differently. They focus more on recent market direction, while statistical gates focus more on volatility. Still,  the true underlying driver of learned gates is undiscovered and may not be any of the general parameters used. Finally, increasing the number of experts does not significantly improve the obtained values. This study also shows that walk-forward hyperparameter selection does not transfer for this problem, so a configuration chosen on validation data can perform worse than the untuned default.
